\documentclass[../../book.tex]{subfiles}

\begin{document}

The protocols of \Cref{section:als20,section:lnp22} prove fairly general statements, but their proofs grow \emph{linearly} with the witness size (note, however, that they almost do not depend on the number of constraints!)
For small statements like a couple of range proofs or commitments' openings this is not a problem: say, LNP22 manages to produce proofs in range 10--20~KB.
For circuits with million gates, however, a linear-size proof is no better than sending the witness itself, while hash-based SNARKs such as Ligero, Aurora, or Brakedown produce proofs whose size grows only polylogarithmically (or as a square root) in the circuit size.

LaBRADOR~\cite{C:BeuSei23} (a \emph{Lattice-Based Recursively Amortized Demonstration Of R1CS}) closes this gap.
It proves an R1CS with $2^{20}$ constraints with a proof of about $58$~KB, which is more than an order of magnitude smaller than the proofs of the hash-based systems above, and it relies only on the Module-SIS assumption (\Cref{def:msis}).
The protocol combines two ideas.
The first one is \emph{amortization}: instead of proving something about $r$ secret vectors separately, the prover folds them into a single random linear combination.
The second one is \emph{recursion}, in the spirit of Bulletproofs (\Cref{section:bulletproofs}): the last message of the protocol is itself a witness for a statement of the same form, so the prover can prove knowledge of it instead of sending it.
The key idea of this section is therefore as follows:

\begin{bookquote}
    Whenever the last message of a proof is a solution to a system of the same type as the original one, the prover can recursively prove knowledge of it instead of sending it.
\end{bookquote}

We should state two remarks before delving into technical specifics, though.
First, LaBRADOR is \emph{not} zero-knowledge.
As we will see, the fact that the protocol no longer needs to be zero-knowledge makes it noticeably simpler than those of the previous sections: there are no masks $\mathbf{y}$ and no rejection sampling.
Second, LaBRADOR proves knowledge of a solution to an R1CS, but the core protocol works with a more convenient relation, the so-called \emph{dot product constraint systems}.
Our plan is therefore as follows:
\begin{enumerate}
    \item In \Cref{section:labrador-relation}, we define the dot product constraint systems.
    \item In \Cref{section:labrador-main}, we build the main protocol step by step, starting from an amortized opening proof.
    \item In \Cref{section:labrador-recursion}, we show how to apply the main protocol recursively.
    \item In \Cref{section:labrador-security}, we discuss the parameters and the security of the protocol.
    \item Finally, in \Cref{section:labrador-r1cs}, we reduce R1CS to dot product constraint systems and give the concrete proof sizes.
\end{enumerate}

Throughout this section, we work over the ring $\mathcal{R}_q = \mathbb{Z}_q[T]/(T^d+1)$, where $q$ is chosen so that $T^d+1$ splits into two irreducible factors modulo $q$, exactly as in LNP22 (\Cref{section:lnp22}).
As before, $f^{\ast} := f(T^{-1})$ denotes the conjugate of $f \in \mathcal{R}_q$, $\mathsf{ct}(f)$ denotes its constant coefficient, and $\|\mathbf{f}\|_2$ denotes the $\ell^2$ norm of the coefficient vector of $\mathbf{f} \in \mathcal{R}_q^n$.
We also write $\vec{\mathbf{f}} \in \mathbb{Z}_q^{nd}$ for the coefficient vector of $\mathbf{f}$ when we want to stress that we treat it as a vector over $\mathbb{Z}_q$.

\subsection{Dot Product Constraint Systems}\label{section:labrador-relation}

Before building the protocol, we need to decide which statements it proves.
The relation must satisfy two requirements.
First, it must be expressive enough to encode R1CS.
Second, the checks that the verifier performs at the end of the protocol must again form a statement of the same type, so that the protocol can be applied recursively.
Both requirements are met by systems of \emph{dot product constraints}.

For two vectors $\mathbf{a}, \mathbf{b} \in \mathcal{R}_q^n$, we write $\langle \mathbf{a}, \mathbf{b} \rangle := \sum_{i \in [n]} a_ib_i \in \mathcal{R}_q$ for their dot product over $\mathcal{R}_q$.
A \textbf{dot product function} of $r$ vector arguments $\mathbf{s}_1, \dots, \mathbf{s}_r \in \mathcal{R}_q^n$ is a quadratic function of the form
\begin{equation*}
    f(\mathbf{s}_1, \dots, \mathbf{s}_r) = \sum_{i,j \in [r]} a_{ij}\langle \mathbf{s}_i, \mathbf{s}_j \rangle + \sum_{i \in [r]} \langle \boldsymbol{\varphi}_i, \mathbf{s}_i \rangle + a_0,
\end{equation*}
where $a_{ij}, a_0 \in \mathcal{R}_q$ and $\boldsymbol{\varphi}_i \in \mathcal{R}_q^n$ are public.
Since $\langle \mathbf{s}_i, \mathbf{s}_j \rangle = \langle \mathbf{s}_j, \mathbf{s}_i \rangle$, we may assume without loss of generality that the matrix $\mathbf{A} := (a_{ij})_{i,j}$ is symmetric.

\begin{remark}
    Note that the class of functions above is a less general version of quadratic forms $\mathbf{s}^{\top}\mathbf{Q}\mathbf{s} + \mathbf{p}^{\top}\mathbf{s} + p_0$ (where $\mathbf{s}=(\mathbf{s}_1,\dots,\mathbf{s}_r) \in \mathcal{R}_q^{nr}$) previously used in \Cref{section:lnp22}.
    We further denote such class of functions as $\mathcal{Q}$.
\end{remark}

Now we define the key relation used in LaBRADOR.

\begin{definition}[Principal Relation]\label{def:labrador-relation}
    Fix the \emph{rank} $n \geq 1$, the \emph{multiplicity} $r \geq 1$, and a norm bound $\beta > 0$.
    A statement consists of two families of dot product functions $\mathcal{F}, \mathcal{F}_{\mathsf{ct}} \subseteq \mathcal{Q}$, and the bound $\beta$; the witness consists of $r$ vectors $\mathbf{s}_1, \dots, \mathbf{s}_r \in \mathcal{R}_q^n$.
    The \textbf{principal relation} of LaBRADOR is
    \begin{equation*}
        \mathcal{R}_{\mathsf{dot}} := \left\{
            \begin{aligned}
                & \mathbbm{x} = (\mathcal{F}, \mathcal{F}_{\mathsf{ct}}, \beta) \\
                & \mathbbm{w} = (\mathbf{s}_1, \dots, \mathbf{s}_r)
            \end{aligned}
        \;\middle|\;
        \begin{aligned}
            & f(\mathbf{s}_1, \dots, \mathbf{s}_r) = 0 \; \text{for all} \; f \in \mathcal{F}, \\
            & \mathsf{ct}\big(g(\mathbf{s}_1, \dots, \mathbf{s}_r)\big) = 0 \; \text{for all} \; g \in \mathcal{F}_{\mathsf{ct}}, \\
            & \textstyle\sum_{i \in [r]} \|\mathbf{s}_i\|_2^2 \leq \beta^2
        \end{aligned}
        \right\}.
    \end{equation*}
\end{definition}

The family $\mathcal{F}$ contains constraints that must vanish as whole polynomials, while the family $\mathcal{F}_{\mathsf{ct}}$ contains constraints of which only the constant coefficient must vanish.
The constraints of $\mathcal{F}_{\mathsf{ct}}$ are exactly those that we have encountered in LNP22: by \Cref{lemma:constant-coefficient}, they express linear and quadratic relations over $\mathbb{Z}_q$ among the coefficients of the witness.

\begin{example}\label{example:labrador-relation}
    Let us express several familiar statements as dot product constraints on a single vector $\mathbf{s} \in \mathcal{R}_q^n$ (so that $r = 1$).
    \begin{enumerate}
        \item \textit{Opening of an Ajtai commitment.} Let $\mathbf{t} = \mathbf{A}\mathbf{s}$ for a public $\mathbf{A} \in \mathcal{R}_q^{\kappa \times n}$ with rows $\mathbf{a}_1^{\top}, \dots, \mathbf{a}_{\kappa}^{\top}$. Then $\mathbf{t} = \mathbf{A}\mathbf{s}$ is the collection of $\kappa$ \emph{linear} functions $f_k(\mathbf{s}) = \langle \mathbf{a}_k, \mathbf{s} \rangle - t_k$ in $\mathcal{F}$.
        \item \textit{A single coefficient.} By \Cref{lemma:constant-coefficient}, the $j$-th coefficient of the $i$-th entry of $\mathbf{s}$ equals $\mathsf{ct}(\langle (T^j\mathbf{e}_i)^{\ast}, \mathbf{s} \rangle)$, where $\mathbf{e}_i$ is the $i$-th unit vector. Hence, a statement such as ``this coefficient equals $5$'' becomes the linear function $g(\mathbf{s}) = \langle (T^j\mathbf{e}_i)^{\ast}, \mathbf{s} \rangle - 5$ in family $\mathcal{F}_{\mathsf{ct}}$.
        \item \textit{Norm.} The squared norm $\|\mathbf{s}\|_2^2$ equals $\mathsf{ct}(\langle \mathbf{s}^{\ast}, \mathbf{s} \rangle)$ (recall Exercise~10 in \Cref{section:lattices-exercises}). This expression is not a dot product function of $\mathbf{s}$ alone, since the conjugate $\mathbf{s}^{\ast}$ appears in it. We will see in \Cref{section:labrador-r1cs} how to overcome this by adding $\mathbf{s}^{\ast}$ to the witness. Notice, however, that the norm bound $\beta$ is already a part of the relation.
    \end{enumerate}
\end{example}

\begin{exercise}\label{exercise:labrador-mlwe-key}
    Let $\mathbf{b} = \mathbf{A}\mathbf{s} + \mathbf{e}$ be a Module-LWE public key (\Cref{def:mlwe}).
    Write the statement ``I know $(\mathbf{s}, \mathbf{e})$ with $\mathbf{b} = \mathbf{A}\mathbf{s} + \mathbf{e}$ and $\|\mathbf{s}\|_2^2 + \|\mathbf{e}\|_2^2 \leq \beta^2$'' as an instance of $\mathcal{R}_{\mathsf{dot}}$ with $r = 1$.
\end{exercise}

In \Cref{section:labrador-r1cs}, we show that an R1CS reduces to $\mathcal{R}_{\mathsf{dot}}$.
For now, we take $\mathcal{R}_{\mathsf{dot}}$ as given and build a proof of knowledge for it.

\subsection{The Main Protocol}\label{section:labrador-main}

We now build a proof of knowledge for $\mathcal{R}_{\mathsf{dot}}$.
The trivial protocol, in which the prover sends $\mathbf{s}_1, \dots, \mathbf{s}_r$, has a proof of $rn$ ring elements.
We improve it in four steps: we first shrink the proof with an amortized opening, and then we gradually add the dot product constraints, the constant-coefficient constraints, and the norm bound.

\subsubsection{Attempt \#1: Amortized Opening}

Recall the opening proof for the BDLOP commitment (\Cref{fig:bdlop-opening}): the prover commits to a mask, receives a challenge $c$, and sends a single response $\mathbf{z} = \mathbf{y} + c\mathbf{r}$ that proves knowledge of the committed randomness.
Suppose now that the prover holds $r$ short vectors $\mathbf{s}_1, \dots, \mathbf{s}_r$ and has committed to each of them with the Ajtai commitment $\mathbf{t}_i = \mathbf{A}\mathbf{s}_i$ for a public $\mathbf{A} \in \mathcal{R}_q^{\kappa \times n}$.
Since we do not need zero-knowledge, the mask $\mathbf{y}$ can be dropped.
Moreover, a \emph{single} response can serve all $r$ commitments: the verifier sends independent challenges $c_1, \dots, c_r$, and the prover replies with the random linear combination
\begin{equation*}
    \mathbf{z} := c_1\mathbf{s}_1 + \dots + c_r\mathbf{s}_r.
\end{equation*}
The verifier checks that $\mathbf{z}$ is short and that $\mathbf{A}\mathbf{z} = c_1\mathbf{t}_1 + \dots + c_r\mathbf{t}_r$.
This technique is called an \textbf{amortized opening}, and the protocol is specified in \Cref{fig:labrador-attempt-1}.

\begin{algoproto}[ht]
    \small
    \centering
    \begin{mdframed}[linecolor=black, linewidth=1pt,
                     skipabove=10pt, skipbelow=10pt,
                     innerleftmargin=5pt, innerrightmargin=5pt]
        \begin{minipage}[t]{0.47\linewidth}
            \colorbox{oc-blue-1}{$\mathcal{P}(\mathbbm{x} = (\mathbf{t}_i)_{i \in [r]}, \mathbbm{w} = (\mathbf{s}_i)_{i \in [r]})$:}
            \begin{algoen}
                \item \keyword{receive} $c_1, \dots, c_r$;
                \item \commentgray{amortized opening} \\ \keyword{send} $\mathbf{z} \gets \sum_{i \in [r]} c_i\mathbf{s}_i$.
            \end{algoen}
        \end{minipage}
        \hfill
        \begin{minipage}[t]{0.47\linewidth}
            \colorbox{oc-blue-1}{$\mathcal{V}(\mathbbm{x} = (\mathbf{t}_i)_{i \in [r]}) \to \{0,1\}$:}
            \begin{algoen}
                \item \keyword{send} $c_1, \dots, c_r \leftarrowS \mathcal{C}$;
                \item \keyword{receive} $\mathbf{z}$;
                \item \keyword{assert} $\|\mathbf{z}\|_2 \leq \gamma$;
                \item \keyword{return} $\llbracket \mathbf{A}\mathbf{z} = \sum_{i \in [r]} c_i\mathbf{t}_i \rrbracket$.
            \end{algoen}
        \end{minipage}
    \end{mdframed}
    \caption{Attempt \#1: amortized opening of $r$ Ajtai commitments $\mathbf{t}_i = \mathbf{A}\mathbf{s}_i$. Here, $\mathcal{C} \subseteq \mathcal{R}_q$ is a set of short challenge polynomials and $\gamma$ is a public norm bound.}
    \label{fig:labrador-attempt-1}
\end{algoproto}

The communication consists of $r\kappa$ ring elements for the commitments and $n$ ring elements for $\mathbf{z}$, instead of $rn$ ring elements for the whole witness.
Since the rank $\kappa$ of the commitment is small (a few dozen), this is a large saving whenever both $r$ and $n$ are large.
The price is that the coefficients of $\mathbf{z}$ are larger than those of the $\mathbf{s}_i$: each coefficient of $\mathbf{z}$ is a sum of $r$ products, so its magnitude grows roughly by a factor of $\sqrt{r\tau}$, where $\tau := \max_{c \in \mathcal{C}}\|c\|_2^2$.

\paragraph{The challenge set}
Similarly to the previous sections, the extractor rewinds the prover and subtracts accepting transcripts, so it needs differences of distinct challenges to be invertible.
However, in contrast to the Lyubashevsky protocol, the extractor obtains a \emph{product} of an extracted vector with a challenge difference, and bounding such products requires one more quantity, the \textbf{operator norm}
\begin{equation*}
    \|c\|_{\mathsf{op}} := \sup_{f \in \mathcal{R}\setminus\{0\}} \frac{\|cf\|_2}{\|f\|_2},
\end{equation*}
which measures how much the multiplication by $c$ can stretch a short element.
LaBRADOR therefore uses a challenge set $\mathcal{C}$ that satisfies three conditions:
\begin{itemize}
    \item the difference $c - c'$ of any two distinct challenges is invertible in $\mathcal{R}_q$;
    \item every challenge is short: $\|c\|_2^2 \leq \tau$;
    \item every challenge has a small operator norm: $\|c\|_{\mathsf{op}} \leq \theta$.
\end{itemize}

\textbf{Concrete choice of challenge set.}
Let $N_i: \mathcal{R} \to \mathbb{N}, c \mapsto \#\{j: |c_j| = i\}$ denote the number of coefficients of $c \in \mathcal{R}$ that have an absolute value of $i$.
The LaBRADOR fixes non-negative integers $k_0 + k_1 + k_2 = d$ and uses the following challenge set:
\begin{equation*}
    \mathcal{C} := \big\{ c \in \mathcal{R} : N_0(c) = k_0, \; N_1(c) = k_1, \; N_2(c) = k_2, \; \|c\|_{\mathsf{op}} \leq \theta \big\}.
\end{equation*}
The scheme uses the following numerical values: $d = 64$, $(k_0, k_1, k_2) = (23, 31, 10)$, and $\theta = 15$.
The ring is chosen so that differences of distinct challenges are invertible~\cite{EC:LyuSei18}, and the condition on the operator norm is enforced by rejection sampling (about $6$ candidates are sampled per challenge on average).

\begin{exercise}\label{exercise:labrador-challenge-count}
    Express $\tau = \|c\|_2^2$ (for $c \in \mathcal{C}$) and $\#\mathcal{C}$ in terms of $d$, $k_0$, $k_1$, $k_2$.

    \commentgray{Answer: $\tau = k_1 + 4k_2$ and $\#\mathcal{C} = \frac{d!}{k_0!\,k_1!\,k_2!} \cdot 2^{k_1 + k_2}$.}
\end{exercise}

\paragraph{Soundness and weak openings}
As in \Cref{remark:relaxed-extraction}, the extractor does not obtain the vectors $\mathbf{s}_i$ themselves.
Instead, it obtains \textbf{weak openings}: vectors $\mathbf{s}_i^{\ast}$ with $\mathbf{A}\mathbf{s}_i^{\ast} = \mathbf{t}_i$ together with a challenge difference $\bar{c} = c - c'$ such that $\bar{c}\mathbf{s}_i^{\ast}$ is short.
The commitment remains binding for weak openings under Module-SIS, but with a norm bound that is larger by a factor proportional to $\theta$.
Indeed, given two weak openings $(\mathbf{s}^{\ast}, \bar{c})$ and $(\mathbf{s}'^{\ast}, \bar{c}')$ of the same commitment, the vector $\bar{c}'(\bar{c}\mathbf{s}^{\ast}) - \bar{c}(\bar{c}'\mathbf{s}'^{\ast})$ lies in the kernel of $\mathbf{A}$, and its norm is bounded using the operator norms of the challenge differences.
This is why the operator norm of the challenges enters the parameters of LaBRADOR (see \Cref{section:labrador-security}).

\subsubsection{Attempt \#2: Proving a Dot Product Constraint}

The amortized opening proves knowledge of \emph{some} short vectors, but not that they satisfy any constraint.
Suppose for now that the statement consists of a single dot product function
\begin{equation*}
    F(\mathbf{s}_1, \dots, \mathbf{s}_r) = \sum_{i,j \in [r]} a_{ij}\langle \mathbf{s}_i, \mathbf{s}_j \rangle + \sum_{i \in [r]} \langle \boldsymbol{\varphi}_i, \mathbf{s}_i \rangle + a_0,
\end{equation*}
and the prover wants to show that $F(\mathbf{s}_1, \dots, \mathbf{s}_r) = 0$.
The idea is similar to the product proof of ALS20 (\Cref{section:als20}): the verifier evaluates the constraint on the amortized opening $\mathbf{z}$, and the result is a polynomial in the challenges whose coefficients depend on the witness.
Namely, since $\mathbf{z} = \sum_i c_i\mathbf{s}_i$, we have
\begin{equation*}
    \langle \mathbf{z}, \mathbf{z} \rangle = \sum_{i,j \in [r]} \textcolor{oc-indigo-9}{g_{ij}}\,c_ic_j, \qquad \sum_{i \in [r]} \langle \boldsymbol{\varphi}_i, \mathbf{z} \rangle c_i = \sum_{i,j \in [r]} \textcolor{oc-teal-9}{h_{ij}}\,c_ic_j,
\end{equation*}
where the \textbf{garbage polynomials} are
\begin{equation*}
    \textcolor{oc-indigo-9}{g_{ij}} := \langle \mathbf{s}_i, \mathbf{s}_j \rangle, \qquad \textcolor{oc-teal-9}{h_{ij}} := \frac{1}{2}\big(\langle \boldsymbol{\varphi}_i, \mathbf{s}_j \rangle + \langle \boldsymbol{\varphi}_j, \mathbf{s}_i \rangle\big).
\end{equation*}
Both matrices $(g_{ij})$ and $(h_{ij})$ are symmetric, and, crucially, they do not depend on the challenges.
Moreover, the constraint itself is a linear function of the garbage:
\begin{equation*}
    F(\mathbf{s}_1, \dots, \mathbf{s}_r) = \sum_{i,j \in [r]} a_{ij}\textcolor{oc-indigo-9}{g_{ij}} + \sum_{i \in [r]} \textcolor{oc-teal-9}{h_{ii}} + a_0.
\end{equation*}
Hence, the prover sends the garbage polynomials together with $\mathbf{z}$, and the verifier checks the three equations above.
The resulting protocol is specified in \Cref{fig:labrador-attempt-2}.

\begin{algoproto}[ht]
    \small
    \centering
    \setlength{\fboxsep}{1.5pt}
    \begin{mdframed}[linecolor=black, linewidth=1pt,
                     skipabove=10pt, skipbelow=10pt,
                     innerleftmargin=5pt, innerrightmargin=5pt]
        \begin{minipage}[t]{0.43\linewidth}
            \colorbox{oc-blue-1}{$\mathcal{P}(\mathbbm{x}, \mathbbm{w} = (\mathbf{s}_i)_{i \in [r]})$:}
            \begin{algoen}
                \item \colorbox{oc-orange-1}{$g_{ij} \gets \langle \mathbf{s}_i, \mathbf{s}_j \rangle$ for $i \leq j$;}
                \item \colorbox{oc-orange-1}{$h_{ij} \gets \frac{1}{2}(\langle \boldsymbol{\varphi}_i, \mathbf{s}_j \rangle + \langle \boldsymbol{\varphi}_j, \mathbf{s}_i \rangle)$} \\ \colorbox{oc-orange-1}{for $i \leq j$;}
                \item \keyword{send} \colorbox{oc-orange-1}{$(g_{ij}, h_{ij})_{i \leq j}$};
                \item \keyword{receive} $c_1, \dots, c_r$;
                \item \keyword{send} $\mathbf{z} \gets \sum_{i \in [r]} c_i\mathbf{s}_i$.
            \end{algoen}
        \end{minipage}
        \hfill
        \begin{minipage}[t]{0.54\linewidth}
            \colorbox{oc-blue-1}{$\mathcal{V}(\mathbbm{x} = ((\mathbf{t}_i)_{i \in [r]}, F)) \to \{0,1\}$:}
            \begin{algoen}
                \item \keyword{receive} $(g_{ij}, h_{ij})_{i \leq j}$;
                \item \keyword{send} $c_1, \dots, c_r \leftarrowS \mathcal{C}$;
                \item \keyword{receive} $\mathbf{z}$;
                \item \keyword{assert} $\|\mathbf{z}\|_2 \leq \gamma$;
                \item \keyword{assert} $\mathbf{A}\mathbf{z} = \sum_{i} c_i\mathbf{t}_i$;
                \item \colorbox{oc-orange-1}{\keyword{assert} $\langle \mathbf{z}, \mathbf{z} \rangle = \sum_{i,j} g_{ij}c_ic_j$;}
                \item \colorbox{oc-orange-1}{\keyword{assert} $\sum_{i} \langle \boldsymbol{\varphi}_i, \mathbf{z} \rangle c_i = \sum_{i,j} h_{ij}c_ic_j$;}
                \item \keyword{return} \colorbox{oc-orange-1}{$\llbracket \sum_{i,j} a_{ij}g_{ij} + \sum_{i} h_{ii} + a_0 = 0 \rrbracket$}.
            \end{algoen}
        \end{minipage}
    \end{mdframed}
    \caption{Attempt \#2: amortized opening together with a proof of a single dot product constraint $F(\mathbf{s}_1, \dots, \mathbf{s}_r) = 0$. Here, $g_{ji} := g_{ij}$ and $h_{ji} := h_{ij}$. Changes compared to \Cref{fig:labrador-attempt-1} are highlighted.}
    \label{fig:labrador-attempt-2}
\end{algoproto}

\begin{exercise}\label{exercise:labrador-garbage}
    Check that the verification equation $\sum_{i} \langle \boldsymbol{\varphi}_i, \mathbf{z} \rangle c_i = \sum_{i,j} h_{ij}c_ic_j$ holds for an honest prover.
\end{exercise}

Why is this sound?
The garbage polynomials are sent \emph{before} the challenges, and the first two new checks compare two quadratic polynomials in the challenges $c_1, \dots, c_r$.
If the prover lies about some $g_{ij}$ or $h_{ij}$, then the two sides of the corresponding check are different polynomials in the challenges, and they agree on random challenges only with a negligible probability.
This is the same argument as in the Schwartz--Zippel lemma, with the only difference that it is applied over $\mathcal{R}_q$.
Therefore, the garbage polynomials are indeed correct, and the last check shows that $F(\mathbf{s}_1, \dots, \mathbf{s}_r) = 0$.

Notice the price of this attempt: the prover sends $r^2 + r$ garbage polynomials, which grows quadratically in $r$.
We will return to this issue in \Cref{section:labrador-recursion}.

\subsubsection{Attempt \#3: Many Constraints}

The statement of $\mathcal{R}_{\mathsf{dot}}$ contains many functions in two families, while Attempt \#2 handles a single function of $\mathcal{F}$.
Both families are reduced to a single function by taking random linear combinations, but the two families apparently require different approaches.

\paragraph{Aggregating $\mathcal{F}$}
If dot product functions $f_1,\dots,f_K \in \mathcal{F}$ must all vanish, then the verifier sends random $\alpha_1, \dots, \alpha_K \leftarrowS \mathcal{R}_q$, and the prover proves that
\begin{equation*}
    F := \sum_{k \in [K]} \alpha_kf_k
\end{equation*}
vanishes.
Since a linear combination of dot product functions is again a dot product function, Attempt \#2 applies directly.
If some $f_k(\mathbf{s}_1, \dots, \mathbf{s}_r) \neq 0$, then the combination vanishes only with small probability over the choice of the $\alpha_k$.

\paragraph{Aggregating $\mathcal{F}_{\mathsf{ct}}$}
For the functions $g_1,\dots,g_L \in \mathcal{F}_{\mathsf{ct}}$, only the constant coefficients vanish, so the combination coefficients must come from $\mathbb{Z}_q$ since multiplying a polynomial with zero constant coefficient by a general element of $\mathcal{R}_q$ may produce a non-zero constant coefficient.
Since a single combination with coefficients from $\mathbb{Z}_q$ has soundness error $1/q$, the verifier sends $K_{\mathsf{ct}} := \lceil \lambda/\log q \rceil$ (for $\lambda=128$ chosen in LaBRADOR) independent combinations with corresponding challenges $\boldsymbol{\psi}_1, \dots, \boldsymbol{\psi}_{K_{\mathsf{ct}}} \leftarrowS \mathbb{Z}_q^L$, and the prover considers
\begin{equation*}
    G_k := \sum_{\ell \in [L]} \psi_{k,\ell}g_{\ell}, \quad k \in [K_{\mathsf{ct}}].
\end{equation*}
For $q \approx 2^{32}$, only $K_{\mathsf{ct}} = 4$ functions are required.

It remains to turn these $K_{\mathsf{ct}}$ constant-coefficient constraints into ordinary ones, and here the lack of zero-knowledge helps.
Write
\begin{equation*}
    G_k(\mathbf{s}_1,\dots,\mathbf{s}_r) = Q_k(\mathbf{s}_1, \dots, \mathbf{s}_r) + b_k,
\end{equation*}
where $Q_k$ collects the quadratic and linear terms.
Recall that LNP22 (\Cref{fig:lnp22-ct}) had to hide the non-constant coefficients of such expressions with a masking polynomial.
LaBRADOR does not need to hide them, so the prover simply sends the polynomial $v_k := Q_k(\mathbf{s}_1, \dots, \mathbf{s}_r)$ in the clear.
The verifier checks that $\mathsf{ct}(v_k) = -\mathsf{ct}(b_k)$, and the claim that $v_k$ was computed correctly is the \emph{ordinary} dot product constraint
\begin{equation*}
    Q_k(\mathbf{s}_1, \dots, \mathbf{s}_r) - v_k = 0,
\end{equation*}
which is added to the family $\mathcal{F}$ before aggregating it.
As a result, the whole statement collapses into a single function $F$, which is proven with Attempt \#2 at the cost of only $K_{\mathsf{ct}}$ additional polynomials $v_k$.

\subsubsection{Attempt \#4: Proving the Norm Bound}

The last missing part of $\mathcal{R}_{\mathsf{dot}}$ is the norm bound $\sum_i \|\mathbf{s}_i\|_2^2 \leq \beta^2$.
The check $\|\mathbf{z}\|_2 \leq \gamma$ in Attempt \#2 is not enough: it only yields weak openings, whose norm may be larger than $\beta$ by a large factor (depending on $\theta$ and $r$).
As we will see in \Cref{section:labrador-r1cs}, the reduction from R1CS needs a norm bound with a small slack, since it uses the bound to rule out wraparound modulo $q$.

The solution is the Johnson--Lindenstrauss projection that we have already met in the approximate range proof of LNP22 (\Cref{lemma:jl-projection}).
Recall that $\mathcal{B}$ denotes the distribution on $\{0, \pm 1\}$ with $\Pr[0] = 1/2$ and $\Pr[\pm 1] = 1/4$.
The verifier sends random matrices $\boldsymbol{\Pi}_1, \dots, \boldsymbol{\Pi}_r \leftarrowS \mathcal{B}^{256 \times nd}$, and the prover sends the projection
\begin{equation*}
    \mathbf{p} := \sum_{i \in [r]} \boldsymbol{\Pi}_i\vec{\mathbf{s}}_i \in \mathbb{Z}_q^{256},
\end{equation*}
where, as before, $\vec{\mathbf{s}}_i \in \mathbb{Z}_q^{nd}$ is the coefficient vector of $\mathbf{s}_i$.
The verifier checks that $\|\mathbf{p}\|_2 \leq \sqrt{128}\beta$.
The following strengthening of \Cref{lemma:jl-projection} justifies this check.

\begin{lemma}[{\cite[Lemma~4.2]{C:BeuSei23}}]\label{lemma:labrador-jl}
    Let $b \leq q/125$ and let $\mathbf{w} \in \mathbb{Z}_q^m$ be any vector with $\|\mathbf{w}\|_2 \geq b$.
    Then,
    \begin{equation*}
        \Pr_{\boldsymbol{\Pi} \leftarrowS \mathcal{B}^{256 \times m}}\Big[\|\boldsymbol{\Pi}\mathbf{w} \bmod q\|_2 < \sqrt{30}\,b\Big] \lesssim 2^{-128}.
    \end{equation*}
\end{lemma}

In other words, a projection of a long vector is long, even after the reduction modulo $q$.
In the honest case, each coordinate of $\boldsymbol{\Pi}\mathbf{w}$ has variance $\|\mathbf{w}\|_2^2/2$, so $\mathbb{E}\big[\|\boldsymbol{\Pi}\mathbf{w}\|_2^2\big] = 128\|\mathbf{w}\|_2^2$.
Applying the lemma to the concatenation $\mathbf{w} = (\vec{\mathbf{s}}_1,\dots,\vec{\mathbf{s}}_r)$ shows the following:
if $\|\mathbf{p}\|_2 \leq \sqrt{128}\beta$, then $\sum_i\|\mathbf{s}_i\|_2^2 \leq \frac{128}{30}\beta^2$ except with probability $2^{-128}$.
Hence, the proof certifies the norm bound up to the slack $\sqrt{128/30} \approx 2.07$.
The honest projection passes the check only with probability about $1/2$, but this affects only the running time of the prover, who simply asks for new projection matrices (or, after the Fiat--Shamir transform, recomputes them).

\begin{exercise}\label{exercise:labrador-jl-mean}
    Let $\mathbf{w} \in \mathbb{Z}^m$ and $\boldsymbol{\pi} \leftarrowS \mathcal{B}^m$. Show that $\mathbb{E}\big[\langle \boldsymbol{\pi}, \mathbf{w} \rangle^2\big] = \frac{1}{2}\|\mathbf{w}\|_2^2$, and deduce that $\mathbb{E}\big[\|\boldsymbol{\Pi}\mathbf{w}\|_2^2\big] = 128\|\mathbf{w}\|_2^2$ for $\boldsymbol{\Pi} \leftarrowS \mathcal{B}^{256 \times m}$.
\end{exercise}

It remains to prove that $\mathbf{p}$ was computed correctly, and this comes almost for free.
Let $\boldsymbol{\pi}_{i,j} \in \mathcal{R}_q^n$ be the vector of polynomials whose coefficient vector is the $j$-th row of $\boldsymbol{\Pi}_i$.
By \Cref{lemma:constant-coefficient}, the $j$-th coordinate of $\mathbf{p}$ is correct if and only if
\begin{equation*}
    \mathsf{ct}\Big(\sum_{i \in [r]} \big\langle \boldsymbol{\pi}_{i,j}^{\ast}, \mathbf{s}_i \big\rangle - p_j\Big) = 0,
\end{equation*}
which is a constant-coefficient constraint.
Hence, the prover simply adds these $256$ functions to the family $\mathcal{F}_{\mathsf{ct}}$, and they are aggregated together with the other functions of $\mathcal{F}_{\mathsf{ct}}$ in Attempt \#3.

\paragraph{The complete protocol}
Putting Attempts \#1--\#4 together gives the base protocol of \Cref{fig:labrador-base}.
Note the order of the messages: the projection must be computed before the aggregation challenges are sent (since the projection constraints are aggregated), and the garbage polynomials $h_{ij}$ can only be computed after the aggregation (since they depend on the aggregated vectors $\boldsymbol{\varphi}_i$).

\begin{algoproto}[ht]
    \footnotesize
    \centering
    \begin{mdframed}[linecolor=black, linewidth=1pt,
                     skipabove=10pt, skipbelow=10pt,
                     innerleftmargin=5pt, innerrightmargin=5pt]
        \begin{minipage}[t]{0.49\linewidth}
            \colorbox{oc-blue-1}{$\mathcal{P}(\mathbbm{x} = (\mathcal{F}, \mathcal{F}_{\mathsf{ct}}, \beta), \mathbbm{w} = (\mathbf{s}_i)_{i \in [r]})$:}
            \begin{algoen}
                \item \commentgray{commit} \\ \keyword{send} $\mathbf{t}_i \gets \mathbf{A}\mathbf{s}_i$ for $i \in [r]$;
                \item \keyword{receive} $\boldsymbol{\Pi}_1, \dots, \boldsymbol{\Pi}_r$;
                \item \commentgray{project} \\ \keyword{send} $\mathbf{p} \gets \sum_{i} \boldsymbol{\Pi}_i\vec{\mathbf{s}}_i$;
                \item \keyword{receive} $\boldsymbol{\psi}_k$ for $k \in [K_{\mathsf{ct}}]$;
                \item \commentgray{aggregate $\mathcal{F}_{\mathsf{ct}}$} \\ \keyword{send} $v_k \gets Q_k(\mathbf{s}_1, \dots, \mathbf{s}_r)$, $k \in [K_{\mathsf{ct}}]$;
                \item \keyword{receive} $\boldsymbol{\alpha}$;
                \item \commentgray{aggregate into $F$} \\ compute $(a_{ij})$, $(\boldsymbol{\varphi}_i)$, $a_0$ of $F$;
                \item \commentgray{garbage terms} \\ \keyword{send} $g_{ij}, h_{ij}$ for $i \leq j$;
                \item \keyword{receive} $c_1, \dots, c_r$;
                \item \commentgray{amortize} \\ \keyword{send} $\mathbf{z} \gets \sum_{i} c_i\mathbf{s}_i$.
            \end{algoen}
        \end{minipage}
        \hfill
        \begin{minipage}[t]{0.49\linewidth}
            \colorbox{oc-blue-1}{$\mathcal{V}(\mathbbm{x} = (\mathcal{F}, \mathcal{F}_{\mathsf{ct}}, \beta)) \to \{0,1\}$:}
            \begin{algoen}
                \item \keyword{receive} $(\mathbf{t}_i)_{i \in [r]}$;
                \item \keyword{send} $\boldsymbol{\Pi}_i \leftarrowS \mathcal{B}^{256 \times nd}$ for $i \in [r]$;
                \item \keyword{receive} $\mathbf{p}$;
                \item \keyword{assert} $\|\mathbf{p}\|_2 \leq \sqrt{128}\beta$;
                \item \keyword{send} $\boldsymbol{\psi}_k \leftarrowS \mathbb{Z}_q^{L+256}$ for $k \in [K_{\mathsf{ct}}]$;
                \item \keyword{receive} $(v_k)_{k \in [K_{\mathsf{ct}}]}$;
                \item \keyword{assert} $\mathsf{ct}(v_k) = -\mathsf{ct}(b_k)$, $k \in [K_{\mathsf{ct}}]$;
                \item \keyword{send} $\boldsymbol{\alpha} \leftarrowS \mathcal{R}_q^{K+K_{\mathsf{ct}}}$;
                \item compute $(a_{ij})$, $(\boldsymbol{\varphi}_i)$, $a_0$ of $F$;
                \item \keyword{receive} $(g_{ij}, h_{ij})_{i \leq j}$;
                \item \keyword{send} $c_1, \dots, c_r \leftarrowS \mathcal{C}$;
                \item \keyword{receive} $\mathbf{z}$;
                \item \keyword{assert} $\|\mathbf{z}\|_2 \leq \gamma$;
                \item \keyword{assert} $\mathbf{A}\mathbf{z} = \sum_{i} c_i\mathbf{t}_i$;
                \item \keyword{assert} $\langle \mathbf{z}, \mathbf{z} \rangle = \sum_{i,j} g_{ij}c_ic_j$;
                \item \keyword{assert} $\sum_{i} \langle \boldsymbol{\varphi}_i, \mathbf{z} \rangle c_i = \sum_{i,j} h_{ij}c_ic_j$;
                \item \keyword{return} $\llbracket \sum_{i,j} a_{ij}g_{ij} + \sum_{i} h_{ii} + a_0 = 0 \rrbracket$.
            \end{algoen}
        \end{minipage}
    \end{mdframed}
    \caption{The base protocol for $\mathcal{R}_{\mathsf{dot}}$ (without recursion). The family $\mathcal{F}_{\mathsf{ct}}$ is extended with the $256$ projection constraints, $\boldsymbol{\psi}_k$ are the coefficients of the $k$-th combination of the extended family, $Q_k$ and $b_k$ are the non-constant part and the constant term of this combination, and $F$ is the combination of $\mathcal{F}$ and the $K_{\mathsf{ct}} = \lceil 128/\log q \rceil$ constraints $Q_k - v_k$ with coefficients $\boldsymbol{\alpha}$.}
    \label{fig:labrador-base}
\end{algoproto}

\subsection{Recursion}\label{section:labrador-recursion}

The base protocol of \Cref{fig:labrador-base} is not succinct yet.
Its proof contains $r\kappa$ commitments, $r^2 + r$ garbage polynomials, and the amortized opening $\mathbf{z}$ of rank $n$.
However, look closely at what the verifier does with the last message: all its checks are dot product constraints.
Indeed, the equation $\mathbf{A}\mathbf{z} = \sum_i c_i\mathbf{t}_i$ is linear in $\mathbf{z}$ and $\mathbf{t}_i$; the equation $\langle \mathbf{z}, \mathbf{z} \rangle = \sum_{i,j} g_{ij}c_ic_j$ is quadratic in $\mathbf{z}$ and linear in the $g_{ij}$; and the remaining two equations are linear in $\mathbf{z}$, $g_{ij}$, and $h_{ij}$.
All coefficients of these constraints (the challenges $c_i$, the matrix $\mathbf{A}$, and the aggregated $a_{ij}$, $\boldsymbol{\varphi}_i$, $a_0$) are known to the verifier.
Therefore, the last message is a witness for a \emph{new} instance of $\mathcal{R}_{\mathsf{dot}}$, and instead of sending it, the prover can prove knowledge of it by running the protocol again.

Before we turn this observation into a protocol, let us make precise what it means to ``prove knowledge of the last message instead of sending it''.
The same idea appears well beyond lattices, so we first look at a more familiar example.

\subsubsection{Warm-up: Proving Knowledge of a Proof}

Suppose we have two proof systems with complementary strengths.
The first one, $(\mathsf{Setup}_{\mathsf{in}}, \mathcal{P}_{\mathsf{in}}, \mathcal{V}_{\mathsf{in}})$, has a fast prover but large proofs; think of a FRI-based STARK, whose proofs take about $100$~KB.
The second one, $(\mathsf{Setup}_{\mathsf{out}}, \mathcal{P}_{\mathsf{out}}, \mathcal{V}_{\mathsf{out}})$, has a slow prover but tiny proofs and constant-time verification; think of Groth16 (\Cref{section:groth}), whose proofs consist of three group elements.
We would like to get the best of both schemes: fast proving and efficient verification.

The idea is to run them one after another, as shown in \Cref{fig:two-level-recursion}.
First, the prover proves the original statement with the fast system and obtains a large proof $\pi_{\mathsf{in}}$.
Then it treats $\pi_{\mathsf{in}}$ as a \emph{witness} for the new relation
\begin{equation*}
    \mathcal{R}_{\mathsf{out}} := \big\{ (\mathbbm{x}_{\mathsf{out}} := (\mathsf{vp}, \mathbbm{x}_{\mathsf{in}}), \mathbbm{w}_{\mathsf{out}} := \pi_{\mathsf{in}}) : \mathcal{V}_{\mathsf{in}}(\mathsf{vp}, \mathbbm{x}_{\mathsf{in}}, \pi_{\mathsf{in}}) = 1 \big\},
\end{equation*}
which says ``I know a proof $\pi_{\mathsf{in}}$ that the inner verifier accepts'', and proves it with Groth16.
The verifier receives only the short proof $\pi_{\mathsf{out}}$ and runs $\mathcal{V}_{\mathsf{out}}$.
The expensive part of the computation (proving the original circuit) is done by the fast prover, while the slow prover handles only the verification circuit of $\mathcal{V}_{\mathsf{in}}$, which is much smaller than the original circuit.
This technique is called \textbf{proof compression}, and it is used in practice to post STARK proofs on-chain.

\begin{figure}[ht]
    \centering
    \begin{tikzpicture}[
        box/.style={draw=oc-indigo-7, thick, fill=oc-indigo-1, rounded corners, align=center, inner sep=4pt, font=\footnotesize},
        dat/.style={font=\footnotesize, align=center},
        arr/.style={-stealth, very thick, oc-gray-7},
    ]
        \node[dat] (in) at (0, 0) {public: $\mathbbm{x}_{\mathsf{in}}$ \\ witness: $\mathbbm{w}_{\mathsf{in}}$};
        \node[box] (p1) at (3.0, 0) {prover $\mathcal{P}_{\mathsf{in}}$ \\ (e.g., STARK)};
        \node[dat] (mid) at (7.2, 0) {public: $(\mathsf{vp}, \mathbbm{x}_{\mathsf{in}})$ \\ witness: $\pi_{\mathsf{in}}$};
        \node[box] (p2) at (7.2, -2.0) {prover $\mathcal{P}_{\mathsf{out}}$ \\ (e.g., Groth16)};
        \node[box, draw=oc-orange-8, fill=oc-orange-1] (v) at (0, -2.0) {verifier $\mathcal{V}_{\mathsf{out}}$};
        \draw[arr] (in) -- (p1);
        \draw[arr] (p1) -- node[above, font=\scriptsize] {$\pi_{\mathsf{in}}$} node[below, font=\scriptsize] {$\approx 100$~KB} (mid);
        \draw[arr] (mid) -- (p2);
        \draw[arr] (p2) -- node[above, font=\scriptsize] {$\pi_{\mathsf{out}}$, a few hundred bytes} (v);
        \node[font=\scriptsize, oc-gray-7] at (3.0, -0.75) {proves $\mathtt{C}(\mathbbm{x}_{\mathsf{in}}, \mathbbm{w}_{\mathsf{in}}) = 0$};
        \node[font=\scriptsize, oc-gray-7] at (7.2, -2.75) {proves $\mathcal{V}_{\mathsf{in}}(\mathsf{vp}, \mathbbm{x}_{\mathsf{in}}, \pi_{\mathsf{in}}) = 1$};
    \end{tikzpicture}
    \caption{Two-level recursion (proof compression). The inner proof system proves the original statement; the outer proof system proves knowledge of an accepting inner proof. The final verifier only runs $\mathcal{V}_{\mathsf{out}}$.}
    \label{fig:two-level-recursion}
\end{figure}

Why is the composed system knowledge sound? We give the following proposition.

\begin{proposition}\label{prop:two-level-soundness}
    Let $\kappa_{\mathsf{in}}$ and $\kappa_{\mathsf{out}}$ be the knowledge errors of inner and outer systems.
    The resultant proving system has a knowledge error of $\kappa_{\mathsf{in}} + \kappa_{\mathsf{out}}$.
\end{proposition}

\begin{proof}
Recall that a proof system is knowledge sound with knowledge error $\kappa$ if there is an extractor $\mathcal{E}$ with the following property: whenever a prover convinces the verifier with probability $\epsilon > \kappa$, the extractor, given access to this prover, outputs a valid witness with probability at least $\epsilon - \kappa$.
Let inner and outer system extractors be $\mathcal{E}_{\mathsf{in}}$ and $\mathcal{E}_{\mathsf{out}}$, respectively.
Now let $\mathcal{A}$ be a prover that convinces $\mathcal{V}_{\mathsf{out}}$ with probability $\epsilon$.
Since the outer system is knowledge sound, its extractor $\mathcal{E}_{\mathsf{out}}$ obtains from $\mathcal{A}$ an accepting inner proof $\pi_{\mathsf{in}}$ with probability at least $\epsilon - \kappa_{\mathsf{out}}$.
But then $\mathcal{E}_{\mathsf{out}}$ (run on $\mathcal{A}$) is itself a successful prover for the inner system, so the inner extractor $\mathcal{E}_{\mathsf{in}}$, run on $\mathcal{E}_{\mathsf{out}}$, obtains a witness $\mathbbm{w}$ with probability at least $\epsilon - (\kappa_{\mathsf{out}} + \kappa_{\mathsf{in}})$.
\end{proof}

\begin{remark}
Note that the argument above has numerous caveats, though:
\begin{enumerate}
    \item the extractors are \emph{nested}: if each extractor runs twice as long as the prover it extracts from, then an extractor for $k$ levels runs $2^k$ times longer, so the argument only works for a small number of levels (at most logarithmic in the security parameter);
    \item if the inner system uses the Fiat-Shamir transform, then the verification circuit contains a concrete hash function, and the security proof requires the (heuristic) assumption that the proof system remains secure when the random oracle is replaced by this hash function.
\end{enumerate}
\end{remark}

\subsubsection{Reductions of Knowledge}

In the example above, the outer system proves a statement that the inner verifier computes from the transcript.
Let us formalize the situation in which an interactive protocol does not decide whether to accept, but instead \emph{reduces} a claim about one relation to a claim about another.
We follow the language of \emph{reductions of knowledge} of Kothapalli and Parno~\cite{C:KotPar23}, as presented in LatticeFold~\cite{EPRINT:BonChe24}.
Throughout, a relation $\mathcal{R}$ consists of triples $(\mathsf{pp}, \mathbbm{x}, \mathbbm{w})$ of public parameters, a statement, and a witness.

\begin{definition}[Reduction of Knowledge]\label{def:reduction-of-knowledge}
    A \textbf{reduction of knowledge} $\Pi$ from a relation $\mathcal{R}_1$ to a relation $\mathcal{R}_2$ consists of an algorithm $\mathsf{Setup}$ and an interactive protocol $\langle \mathcal{P}, \mathcal{V} \rangle$:
    \begin{itemize}
        \item $\mathsf{Setup}(1^{\lambda}) \to \mathsf{pp}$ outputs the public parameters;
        \item $\langle \mathcal{P}(\mathsf{pp}, \mathbbm{x}_1, \mathbbm{w}_1), \mathcal{V}(\mathsf{pp}, \mathbbm{x}_1) \rangle \to (\mathbbm{x}_2, \mathbbm{w}_2)$: at the end of the interaction, the verifier outputs a new statement $\mathbbm{x}_2$ (or $\bot$ to reject), and the prover outputs a new witness $\mathbbm{w}_2$.
    \end{itemize}
\end{definition}

Now we give the security notions that we require from the reduction of knowledge.

\begin{definition}
    The reduction of knowledge $\langle \mathcal{P}, \mathcal{V} \rangle: \mathcal{R}_1 \to \mathcal{R}_2$ may satisfy the following properties.
    \begin{itemize}
        \item \textbf{Completeness.} For all $\mathsf{pp} \gets \mathsf{Setup}(1^{\lambda})$ and all $(\mathsf{pp}, \mathbbm{x}_1, \mathbbm{w}_1) \in \mathcal{R}_1$, the output $(\mathbbm{x}_2, \mathbbm{w}_2)$ of the honest interaction satisfies $(\mathsf{pp}, \mathbbm{x}_2, \mathbbm{w}_2) \in \mathcal{R}_2$.
        \item \textbf{Knowledge soundness} (with knowledge error $\kappa$). There exists an expected polynomial-time extractor $\mathcal{E}$ such that the following holds for every adversary $(\mathcal{A}, \mathcal{P}^{\ast})$ running in expected polynomial time. If
        \begin{equation*}
            \Pr\left[(\mathsf{pp}, \mathbbm{x}_2, \mathbbm{w}_2) \in \mathcal{R}_2 \;\middle|\;
            \begin{aligned}
                & \mathsf{pp} \gets \mathsf{Setup}(1^{\lambda}), \\
                & (\mathbbm{x}_1, \mathsf{st}) \gets \mathcal{A}(\mathsf{pp}), \\
                & (\mathbbm{x}_2, \mathbbm{w}_2) \gets \langle \mathcal{P}^{\ast}(\mathsf{st}), \mathcal{V}(\mathsf{pp}, \mathbbm{x}_1) \rangle
            \end{aligned}
            \right] = \epsilon > \kappa,
        \end{equation*}
        then $\mathcal{E}^{\mathcal{A}, \mathcal{P}^{\ast}}(\mathsf{pp}, \mathbbm{x}_1)$ outputs $\mathbbm{w}_1$ with $(\mathsf{pp}, \mathbbm{x}_1, \mathbbm{w}_1) \in \mathcal{R}_1$ with probability at least $\epsilon - \kappa$.
        \item \textbf{Public reducibility.} There is a deterministic poly-time algorithm $\mathsf{f}$ that computes the output statement from the transcript: $\mathbbm{x}_2 = \mathsf{f}(\mathsf{pp}, \mathbbm{x}_1, \mathsf{tr})$.
    \end{itemize}
\end{definition}

In words, whoever can produce a witness for the new statement $\mathbbm{x}_2$ must have known a witness for the original statement $\mathbbm{x}_1$.
Public reducibility says that $\mathbbm{x}_2$ is determined by the transcript, so that a later protocol (or a Fiat--Shamir verifier) can recompute it.

\begin{example}\label{example:reductions-of-knowledge}
    In fact, we can already name at least three examples of reductions of knowledge:
    \begin{enumerate}
        \item An argument of knowledge for $\mathcal{R}$ with knowledge error $\kappa$ is exactly a reduction of knowledge from $\mathcal{R}$ to the trivial relation $\mathcal{R}_{\top} := \{(\mathsf{pp}, \top, \bot)\}$ with knowledge error $\kappa$, whose only statement $\top$ means ``accept'': the verifier outputs $\top$ if it accepts and $\bot$ otherwise, and the prover outputs the empty witness $\bot$.
        \item The proof compression from \Cref{fig:two-level-recursion} is a sequence of two reductions: given that the inner proving system proves a relation $\mathcal{R}_{\mathsf{in}}$ while the outer proves $\mathcal{R}_{\mathsf{out}}$, the reduction is
        \begin{equation*}
            \mathcal{R}_{\mathsf{in}} \to \mathcal{R}_{\mathsf{out}} \to \mathcal{R}_{\top}
        \end{equation*}
        with a composite knowledge error of $\kappa_{\mathsf{in}} + \kappa_{\mathsf{out}}$, as was previously shown in \Cref{prop:two-level-soundness}.
        \item Every relation $\mathcal{R}$ has the trivial argument of knowledge $\Pi_{\mathsf{send}}$, in which the prover sends $\mathbbm{w}$ and the verifier outputs $\top$ if and only if $(\mathsf{pp}, \mathbbm{x}, \mathbbm{w}) \in \mathcal{R}$. 
        It is a reduction of knowledge from $\mathcal{R}$ to $\mathcal{R}_{\top}$ with knowledge error $0$ (the extractor simply reads $\mathbbm{w}$ from the transcript), but it is not succinct.
    \end{enumerate}
\end{example}

The following theorem is what makes reductions of knowledge so convenient: they compose like functions.

\begin{theorem}[Composition~{\cite[Theorems~5 and~6]{C:KotPar23}}, {\cite[Lemma~3.7]{C:BeuSei23}}]\label{thm:rok-composition}
    Let $\Pi$ be a reduction of knowledge from $\mathcal{R}$ to $\mathcal{R}'$ with knowledge error $\kappa$.
    \begin{enumerate}
        \item[\textit{(i)}] \textbf{Sequential composition.} If $\Pi'$ is a reduction of knowledge from $\mathcal{R}'$ to $\mathcal{R}''$ with knowledge error $\kappa'$, then the protocol $\Pi' \circ \Pi$, which runs $\Pi$ and then runs $\Pi'$ on its output, is a reduction of knowledge from $\mathcal{R}$ to $\mathcal{R}''$ with knowledge error $\kappa + \kappa'$.
        \item[\textit{(ii)}] \textbf{Parallel composition.} If $\widetilde{\Pi}$ is a reduction of knowledge from $\widetilde{\mathcal{R}}$ to $\widetilde{\mathcal{R}}'$, then running $\Pi$ and $\widetilde{\Pi}$ in parallel gives a reduction of knowledge $\Pi \times \widetilde{\Pi}$ from $\mathcal{R} \times \widetilde{\mathcal{R}}$ to $\mathcal{R}' \times \widetilde{\mathcal{R}}'$.
    \end{enumerate}
\end{theorem}

The proof of sequential composition is exactly the nested extraction from the warm-up: the extractor of $\Pi'$ turns a successful prover for $\Pi' \circ \Pi$ into a successful prover for $\Pi$, and the extractor of $\Pi$ is then run on it.
The technical work lies in keeping the running time expected polynomial, and this is why \Cref{def:reduction-of-knowledge} quantifies over expected polynomial-time adversaries.

\begin{remark}[Relaxed relations]\label{remark:rok-relaxed}
    Lattice-based reductions are usually complete for one relation but knowledge sound only for a \emph{relaxed} one, as we have already seen in \Cref{remark:relaxed-extraction}.
    For example, the main protocol of LaBRADOR is complete for witnesses of norm at most $\beta$, but the extractor only guarantees norm at most $\sqrt{128/30}\,\beta$.
    When composing such reductions, the relaxed relation that $\Pi$ is knowledge sound for must be handled by $\Pi'$.
    In LaBRADOR, this is taken into account by choosing the Module-SIS parameters of each iteration with the norm slack of the next iterations in mind.
\end{remark}

\subsubsection{LaBRADOR as a Chain of Reductions}

We can now describe the structure of LaBRADOR precisely.
Each iteration of the main protocol is a reduction of knowledge
$\Pi_{\mathsf{main}}: \mathcal{R}_{\mathsf{dot}} \to \mathcal{R}_{\mathsf{dot}}'$
from the principal relation with one set of parameters to the same relation with new parameters (we construct it below).
The reduction from R1CS in \Cref{section:labrador-r1cs} is a reduction of knowledge $\Pi_{\mathsf{R1CS}}: \mathcal{R}_{\mathsf{R1CS}} \to \mathcal{R}_{\mathsf{dot}}$.
The complete protocol is
\begin{equation*}
    \Pi_{\mathsf{LaBRADOR}} := \Pi_{\mathsf{send}} \circ \underbrace{\Pi_{\mathsf{main}} \circ \dots \circ \Pi_{\mathsf{main}}}_{\ell \; \text{times}} \circ \Pi_{\mathsf{R1CS}},
\end{equation*}
and by \Cref{thm:rok-composition}, it is an argument of knowledge for R1CS with knowledge error $\kappa_{\mathsf{R1CS}} + \ell \cdot \kappa_{\mathsf{main}}$, which is still negligible.
Parallel composition is used as well: the reductions for binary R1CS and for R1CS modulo $2^d+1$ are run in parallel, and their outputs are merged into a single instance of $\mathcal{R}_{\mathsf{dot}}$.

Note that the same idea of composing multiple reductions underlies folding schemes such as Nova~\cite{C:KotSetTzi22} and LatticeFold~\cite{EPRINT:BonChe24}, which are reductions of knowledge from two instances of a relation to one.
Moreover, both caveats of the warm-up disappear: the number of iterations is $O(\log\log N)$, and the Fiat-Shamir transform is applied once to the whole composed public-coin protocol, rather than to each level separately.

\subsubsection{Making the Main Protocol Recursive}

To turn the base protocol into a reduction of knowledge from $\mathcal{R}_{\mathsf{dot}}$ to itself, we need to address three problems:
\begin{enumerate}
    \item[\textit{Issue \#1}] the commitments $\mathbf{t}_i$ and the garbage polynomials are still sent in the clear, and their total size grows with $r$;
    \item[\textit{Issue \#2}] the coefficients of $\mathbf{z}$ are larger than those of the $\mathbf{s}_i$ by a factor of about $\sqrt{r\tau}$, so they would blow up after a few iterations;
    \item[\textit{Issue \#3}] the new witness consists of a single long vector $\mathbf{z}$ and several short vectors, while the protocol would benefit most from a witness consisting of many vectors of the same rank.
\end{enumerate}

\paragraph{Outer commitments}
The \textit{Issue \#1} is solved by moving the commitments and the garbage polynomials into the new witness.
The verifier does not need to \emph{see} them, but they must be fixed before the challenges that depend on them.
Hence, instead of sending $\mathbf{t}_1, \dots, \mathbf{t}_r$, the prover sends a single Ajtai commitment to all of them, and similarly for the garbage polynomials.
We call the $\mathbf{t}_i$ the \textbf{inner commitments} and the new commitments the \textbf{outer commitments}.

One subtlety is that the Ajtai commitment requires a short message (see \Cref{section:sis}), while the coefficients of the $\mathbf{t}_i$ are arbitrary elements of $\mathbb{Z}_q$.
Therefore, the prover first decomposes each $\mathbf{t}_i$ in a small base (call it $\delta \in \mathbb{N}$):
\begin{equation*}
    \mathbf{t}_i = \sum_{k \in [t]} \mathbf{t}_{i,k}\delta^k, \quad \|\mathbf{t}_{i,k}\|_{\infty} \leq \frac{1}{2}\delta,
\end{equation*}
where $t = \lceil \log q/\log \delta \rceil$, and similarly decomposes the garbage polynomials $h_{ij}$ in base $\delta$ and the (already relatively short) garbage polynomials $g_{ij}$ in some base $\delta'$.
Denote by $\mathbf{t}$, $\mathbf{g}$, and $\mathbf{h}$ the vectors collecting all the decomposition parts of the inner commitments and of the garbage polynomials.
The prover sends two outer commitments,
\begin{equation*}
    \mathbf{u}' := \mathbf{B}\mathbf{t} + \mathbf{C}\mathbf{g} \in \mathcal{R}_q^{\kappa'}, \qquad \mathbf{u}'' := \mathbf{D}\mathbf{h} \in \mathcal{R}_q^{\kappa''},
\end{equation*}
for public matrices $\mathbf{B}$, $\mathbf{C}$, $\mathbf{D}$.
There are two outer commitments, since they are sent at different times: the inner commitments and the $g_{ij} = \langle \mathbf{s}_i, \mathbf{s}_j \rangle$ depend only on the witness and are committed to in the first message, while the $h_{ij}$ depend on the aggregated vectors $\boldsymbol{\varphi}_i$ and can be computed only after the aggregation challenges.
The equations $\mathbf{u}' = \mathbf{B}\mathbf{t} + \mathbf{C}\mathbf{g}$ and $\mathbf{u}'' = \mathbf{D}\mathbf{h}$ are linear, so they simply become $\kappa' + \kappa''$ additional linear dot product constraints of the next instance.

\paragraph{Decomposing the amortized opening}
The \textit{Issue \#2} is solved in the same way: the prover decomposes $\mathbf{z}$ in a base $\delta'' \in \mathbb{N}$ into two parts,
\begin{equation*}
    \mathbf{z} = \mathbf{z}_{\mathsf{low}} + \delta''\mathbf{z}_{\mathsf{high}}, \quad \|\mathbf{z}_{\mathsf{low}}\|_{\infty} \leq \frac{1}{2}\delta'',
\end{equation*}
and the quadratic check becomes
\begin{equation*}
    \langle \mathbf{z}_{\mathsf{low}}, \mathbf{z}_{\mathsf{low}} \rangle + 2\delta''\langle \mathbf{z}_{\mathsf{low}}, \mathbf{z}_{\mathsf{high}} \rangle + \delta''^2\langle \mathbf{z}_{\mathsf{high}}, \mathbf{z}_{\mathsf{high}} \rangle = \sum_{i,j} g_{ij}c_ic_j,
\end{equation*}
which is still a dot product constraint.
The base $\delta''$ is chosen so that both parts have coefficients of about the same size, and the bases $\delta$, $\delta'$ are chosen so that $\delta \approx \delta' \approx \delta''$, so that all parts of the new witness are about equally wide.

\paragraph{Reshaping the witness}
To solve the \textit{Issue \#3}, the prover cuts the new witness into pieces of the same rank.
Namely, let $\mathbf{o} := \mathbf{t}\,\|\,\mathbf{g}\,\|\,\mathbf{h} \in \mathcal{R}_q^m$ be the concatenation of all openings of the outer commitments.
The prover splits each of $\mathbf{z}_{\mathsf{low}}$ and $\mathbf{z}_{\mathsf{high}}$ into $\nu$ pieces and the vector $\mathbf{o}$ into $\mu$ pieces, and pads all pieces with zeros to the common rank $n' := \max\{\lceil n/\nu \rceil, \lceil m/\mu \rceil\}$.
This gives a new witness $\mathbf{s}_1', \dots, \mathbf{s}_{r'}'$ with multiplicity $r' = 2\nu + \mu$ and rank $n'$, where $\nu$ and $\mu$ are chosen so that $n/\nu \approx m/\mu$ to avoid excessive padding.
Every constraint above is still a dot product constraint in the pieces $\mathbf{s}_i'$, and the three norm checks (on $\mathbf{z}$, on $\mathbf{t}$ and $\mathbf{g}$, and on $\mathbf{h}$) are combined into a single one:
\begin{equation*}
    \sum_{i \in [r']} \|\mathbf{s}_i'\|_2^2 = \|\mathbf{z}_{\mathsf{low}}\|_2^2 + \|\mathbf{z}_{\mathsf{high}}\|_2^2 + \|\mathbf{o}\|_2^2 \leq \beta'^2.
\end{equation*}
The new instance of $\mathcal{R}_{\mathsf{dot}}$ consists of $\kappa + \kappa' + \kappa'' + 3$ constraints of the family $\mathcal{F}$, an empty family $\mathcal{F}_{\mathsf{ct}}$, and the norm bound $\beta'$.
The whole process is illustrated in \Cref{fig:labrador-iteration}, and the complete iteration of the main protocol is specified in \Cref{fig:labrador-main}.

\begin{figure}[ht]
    \centering
    \begin{tikzpicture}[
        wit/.style={draw=oc-indigo-7, thick, fill=oc-indigo-1, minimum height=0.45cm, inner sep=0pt},
        aux/.style={draw=oc-teal-7, thick, fill=oc-teal-1, minimum height=0.45cm, inner sep=0pt},
        prf/.style={draw=oc-orange-8, thick, fill=oc-orange-1, rounded corners, inner sep=3pt, font=\footnotesize},
        lab/.style={font=\footnotesize, anchor=east},
        arr/.style={-stealth, very thick, oc-gray-7},
    ]
        % old witness: r = 6 long vectors
        \node[lab] at (-0.2, 0) {$(\mathbf{s}_i)_{i \in [r]}$};
        \foreach \i in {0,...,5} {
            \node[wit, minimum width=0.9cm] at (0.45+0.95*\i, 0) {};
        }
        \node[font=\scriptsize, oc-gray-7] at (2.85, 0.45) {$r$ vectors of rank $n$};
        % folded witness
        \draw[arr] (2.85,-0.35) -- node[right, font=\footnotesize, black] {fold and commit} (2.85,-1.05);
        \node[lab] at (-0.2, -1.4) {$\mathbf{z}, \mathbf{o}$};
        \node[wit, minimum width=0.9cm] (z) at (0.45, -1.4) {$\scriptstyle\mathbf{z}$};
        \node[aux, minimum width=1.6cm] (o) at (2.0, -1.4) {$\scriptstyle\mathbf{t}\,\|\,\mathbf{g}\,\|\,\mathbf{h}$};
        \node[prf] (pi) at (5.3, -1.4) {$\pi = (\mathbf{u}', \mathbf{p}, (v_k)_k, \mathbf{u}'')$};
        \draw[-stealth, thick, oc-orange-8] (pi.east) -- ++(0.6,0) node[right, font=\footnotesize, black] {$\mathcal{V}$};
        % decomposed & reshaped witness
        \draw[arr] (2.85,-1.75) -- node[right, font=\footnotesize, black] {decompose and reshape} (2.85,-2.45);
        \node[lab] at (-0.2, -2.8) {$(\mathbf{s}_i')_{i \in [r']}$};
        \foreach \i in {0,...,3} {
            \node[wit, minimum width=0.4cm] at (0.2+0.45*\i, -2.8) {};
        }
        \foreach \i in {0,...,3} {
            \node[aux, minimum width=0.4cm] at (2.05+0.45*\i, -2.8) {};
        }
        \node[font=\scriptsize, oc-gray-7] at (1.9, -3.25) {$r' = 2\nu + \mu$ vectors of rank $n' \approx n/\nu$};
    \end{tikzpicture}
    \caption{One iteration of the main protocol. The $r$ witness vectors are folded into the amortized opening $\mathbf{z}$ of rank $n$, and the openings $\mathbf{o} = \mathbf{t}\,\|\,\mathbf{g}\,\|\,\mathbf{h}$ of the outer commitments are appended to it. Only the short proof $\pi$ is sent to the verifier. The new witness is obtained by decomposing $\mathbf{z}$ and cutting everything into $r'$ pieces of rank $n'$.}
    \label{fig:labrador-iteration}
\end{figure}

\begin{algoproto}[ht]
    \footnotesize
    \centering
    \setlength{\fboxsep}{1.5pt}
    \begin{mdframed}[linecolor=black, linewidth=1pt,
                     skipabove=10pt, skipbelow=10pt,
                     innerleftmargin=5pt, innerrightmargin=5pt]
        \begin{minipage}[t]{0.52\linewidth}
            \colorbox{oc-blue-1}{$\mathcal{P}(\mathbbm{x} = (\mathcal{F}, \mathcal{F}_{\mathsf{ct}}, \beta), \mathbbm{w} = (\mathbf{s}_i)_{i \in [r]})$:}
            \begin{algoen}
                \item $\mathbf{t}_i \gets \mathbf{A}\mathbf{s}_i$ for $i \in [r]$;
                \item $g_{ij} \gets \langle \mathbf{s}_i, \mathbf{s}_j \rangle$ for $i \leq j$;
                \item \colorbox{oc-orange-1}{decompose the $\mathbf{t}_i$ and $g_{ij}$ into $\mathbf{t}$ and $\mathbf{g}$;}
                \item \keyword{send} \colorbox{oc-orange-1}{$\mathbf{u}' \gets \mathbf{B}\mathbf{t} + \mathbf{C}\mathbf{g}$};
                \item \keyword{receive} $\boldsymbol{\Pi}_1, \dots, \boldsymbol{\Pi}_r$;
                \item \keyword{send} $\mathbf{p} \gets \sum_{i} \boldsymbol{\Pi}_i\vec{\mathbf{s}}_i$;
                \item \keyword{receive} $\boldsymbol{\psi}_k$ for $k \in [K_{\mathsf{ct}}]$;
                \item \keyword{send} $v_k \gets Q_k(\mathbf{s}_1, \dots, \mathbf{s}_r)$ for $k \in [K_{\mathsf{ct}}]$;
                \item \keyword{receive} $\boldsymbol{\alpha}$;
                \item compute $(a_{ij})$, $(\boldsymbol{\varphi}_i)$, $a_0$ of $F$;
                \item $h_{ij} \gets \frac{1}{2}(\langle \boldsymbol{\varphi}_i, \mathbf{s}_j \rangle + \langle \boldsymbol{\varphi}_j, \mathbf{s}_i \rangle)$ for $i \leq j$;
                \item \colorbox{oc-orange-1}{decompose the $h_{ij}$ into $\mathbf{h}$;}
                \item \keyword{send} \colorbox{oc-orange-1}{$\mathbf{u}'' \gets \mathbf{D}\mathbf{h}$};
                \item \keyword{receive} $c_1, \dots, c_r$;
                \item $\mathbf{z} \gets \sum_{i} c_i\mathbf{s}_i$;
                \item \colorbox{oc-orange-1}{decompose $\mathbf{z} = \mathbf{z}_{\mathsf{low}} + \delta\mathbf{z}_{\mathsf{high}}$;}
                \item \colorbox{oc-orange-1}{\keyword{return} $\mathbbm{w}' \gets (\mathbf{s}_i')_{i \in [r']}$} \\ \colorbox{oc-orange-1}{where $(\mathbf{s}_i')_i \gets \mathsf{Reshape}(\mathbf{z}_{\mathsf{low}}, \mathbf{z}_{\mathsf{high}}, \mathbf{t}, \mathbf{g}, \mathbf{h})$}.
            \end{algoen}
        \end{minipage}
        \hfill
        \begin{minipage}[t]{0.45\linewidth}
            \colorbox{oc-blue-1}{$\mathcal{V}(\mathbbm{x} = (\mathcal{F}, \mathcal{F}_{\mathsf{ct}}, \beta))$:}
            \begin{algoen}
                \item \keyword{receive} \colorbox{oc-orange-1}{$\mathbf{u}'$};
                \item \keyword{send} $\boldsymbol{\Pi}_i \leftarrowS \mathcal{B}^{256 \times nd}$ for $i \in [r]$;
                \item \keyword{receive} $\mathbf{p}$;
                \item \keyword{assert} $\|\mathbf{p}\|_2 \leq \sqrt{128}\beta$;
                \item \keyword{send} $\boldsymbol{\psi}_k \leftarrowS \mathbb{Z}_q^{L+256}$, $k \in [K_{\mathsf{ct}}]$;
                \item \keyword{receive} $(v_k)_{k \in [K_{\mathsf{ct}}]}$;
                \item \keyword{assert} $\mathsf{ct}(v_k) = -\mathsf{ct}(b_k)$, $k \in [K_{\mathsf{ct}}]$;
                \item \keyword{send} $\boldsymbol{\alpha} \leftarrowS \mathcal{R}_q^{K+K_{\mathsf{ct}}}$;
                \item compute $(a_{ij})$, $(\boldsymbol{\varphi}_i)$, $a_0$ of $F$;
                \item \keyword{receive} \colorbox{oc-orange-1}{$\mathbf{u}''$};
                \item \keyword{send} $c_1, \dots, c_r \leftarrowS \mathcal{C}$;
                \item \commentgray{new statement} \\ \colorbox{oc-orange-1}{\keyword{return} $\mathbbm{x}' \gets (\mathcal{F}', \varnothing, \beta')$}.
            \end{algoen}
        \end{minipage}
    \end{mdframed}
    \caption{One iteration of the main protocol of LaBRADOR~\cite[Fig.~2]{C:BeuSei23}. Instead of sending its last message, the prover outputs the new witness $(\mathbf{s}_i')_{i \in [r']}$, and the verifier outputs the new statement: the family $\mathcal{F}'$ consists of the $\kappa + \kappa' + \kappa'' + 3$ constraints $\mathbf{A}\mathbf{z} = \sum_i c_i\mathbf{t}_i$, $\mathbf{u}' = \mathbf{B}\mathbf{t} + \mathbf{C}\mathbf{g}$, $\mathbf{u}'' = \mathbf{D}\mathbf{h}$, and the three garbage equations of \Cref{fig:labrador-base}, all rewritten in terms of the new witness. Constant-coefficient constraints are not needed. We highlight changes from \Cref{fig:labrador-base}.}
    \label{fig:labrador-main}
\end{algoproto}

\paragraph{How fast does the witness shrink?}
Let us count the size of the witness in ring elements, ignoring the small constants.
Before the iteration, the witness consists of $N = rn$ elements.
After the iteration, it consists of about $2n$ elements of $\mathbf{z}_{\mathsf{low}}$ and $\mathbf{z}_{\mathsf{high}}$ and about $r^2$ elements of $\mathbf{o}$ (dominated by the garbage polynomials).
Since the prover is free to cut its witness into any number of pieces, it chooses $r$ to minimize $2N/r + r^2$, which gives $r \approx N^{1/3}$ and a new witness of size about $N^{2/3}$ (up to a constant factor).
Hence, the witness size goes from $N$ to $N^{2/3}$, then to $N^{4/9}$, and so on, and only $O(\log \log N)$ iterations are needed until the witness becomes small.
In practice, six or seven iterations suffice for all relevant circuit sizes.
After the last iteration, the prover sends the remaining witness in the clear, and the verifier checks it directly.

\begin{exercise}\label{exercise:labrador-shrink}
    Suppose the witness has $N = 2^{21}$ ring elements. Demonstrate that the multiplicity $r \approx N^{1/3}$ minimizes $2N/r + r^2$, and compute the size of the new witness.
\end{exercise}

Meanwhile, each iteration contributes only the short proof $\pi$ to the communication: the outer commitments $\mathbf{u}', \mathbf{u}''$ ($\kappa' + \kappa''$ ring elements), the projection $\mathbf{p}$ ($256$ small integers), and the $K_{\mathsf{ct}}$ polynomials $v_k$.
This is a few kilobytes per iteration, independently of the witness size.
The complete recursive protocol is summarized in \Cref{fig:labrador-recursive}.

\begin{algoproto}[H]
    \small
    \centering
    \begin{mdframed}[linecolor=black, linewidth=1pt,
                     skipabove=10pt, skipbelow=10pt,
                     innerleftmargin=5pt, innerrightmargin=5pt]
        \begin{minipage}[t]{0.48\linewidth}
            \colorbox{oc-blue-1}{$\mathcal{P}(\mathbbm{x}_0, \mathbbm{w}_0)$:}
            \begin{algoen}
                \item \keyword{for} $i \in [\ell]$:
                \item \hspace{5px} run \Cref{fig:labrador-main} on $(\mathbbm{x}_{i}, \mathbbm{w}_{i})$;
                \item \hspace{5px} $\mathbbm{w}_{i+1} \gets$ its output;
                \item \hspace{5px} compute $\mathbbm{x}_{i+1}$ as the verifier does;
                \item \commentgray{last iteration} \\ \keyword{send} $\mathbbm{w}_{\ell}$.
            \end{algoen}
        \end{minipage}
        \hfill
        \begin{minipage}[t]{0.48\linewidth}
            \colorbox{oc-blue-1}{$\mathcal{V}(\mathbbm{x}_0) \to \{0,1\}$:}
            \begin{algoen}
                \item \keyword{for} $i \in [\ell]$:
                \item \hspace{5px} run \Cref{fig:labrador-main} on $\mathbbm{x}_{i}$;
                \item \hspace{5px} $\mathbbm{x}_{i+1} \gets$ its output;
                \item \keyword{receive} $\mathbbm{w}_{\ell}$;
                \item \keyword{return} $\llbracket (\mathbbm{x}_{\ell}, \mathbbm{w}_{\ell}) \in \mathcal{R}_{\mathsf{dot}} \rrbracket$.
            \end{algoen}
        \end{minipage}
    \end{mdframed}
    \caption{The recursive LaBRADOR protocol with $\ell$ iterations. If any check inside an iteration fails, the verifier rejects. The proof consists of the short messages $\pi$ of all iterations and the final witness $\mathbbm{w}_{\ell}$.}
    \label{fig:labrador-recursive}
\end{algoproto}

\subsection{Parameters and Security}\label{section:labrador-security}

We do not analyze the extractor of LaBRADOR in detail, but two aspects of the security analysis affect the choice of parameters, so we discuss them here: the norm bounds and the hardness assumptions they require.

\paragraph{Norm bounds}
The honest prover must pass the consolidated norm check $\sum_i \|\mathbf{s}_i'\|_2^2 \leq \beta'^2$, so the bound $\beta'$ is chosen according to the expected norms of the honest vectors.
LaBRADOR estimates these norms heuristically by modelling all coefficients as independent Gaussians, which matches experiments well.
Suppose the coefficients of the $\mathbf{s}_i$ have standard deviation $\sigma = \beta/\sqrt{rnd}$.
Each coefficient of $\mathbf{z} = \sum_i c_i\mathbf{s}_i$ is then a sum of about $rd$ such coefficients multiplied by challenge coefficients, so it has standard deviation about $\sigma\sqrt{r\tau}$, and hence
\begin{equation*}
    \|\mathbf{z}\|_2 \approx \sigma\sqrt{r\tau}\cdot\sqrt{nd} = \beta\sqrt{\tau} =: \gamma.
\end{equation*}
After the decomposition $\mathbf{z} = \mathbf{z}_{\mathsf{low}} + \delta''\mathbf{z}_{\mathsf{high}}$, the low part is essentially uniform modulo $\delta''$ (with standard deviation $\delta''/\sqrt{12}$), while the high part has standard deviation $\sigma\sqrt{r\tau}/\delta''$.
Both are equal for
\begin{equation*}
    \delta'' \approx \sqrt{\sqrt{12r\tau}\cdot \sigma},
\end{equation*}
and the bases $\delta$, $\delta'$ of the outer commitments are then chosen close to $\delta''$, so that the parts of $\mathbf{o}$ are as wide as the parts of $\mathbf{z}$.
The bounds $\gamma'$ and $\gamma''$ on the openings of $\mathbf{u}'$ and $\mathbf{u}''$ are obtained in the same way, and the new bound is $\beta'^2 = \frac{2}{\delta''^2}\gamma^2 + \gamma'^2 + \gamma''^2$.
If an honest execution happens to produce longer vectors than predicted, the prover simply restarts; this affects the running time, but not soundness.

\iffalse
\paragraph{Security}
As discussed in \Cref{section:labrador-main}, the extractor obtains weak openings of the inner commitments, so the Module-SIS assumption for the matrix $\mathbf{A}$ is needed with a norm bound that involves the operator norm $\theta$ of the challenges.
The outer commitments, in contrast, are opened directly in the last message, so they only need Module-SIS with twice the norm of their openings.
Finally, the Johnson-Lindenstrauss argument of \Cref{lemma:labrador-jl} requires the certified norm to be small compared to $q$.
We state the resulting theorem informally.

\begin{theorem}[Informal, {\cite[Theorem~5.1]{C:BeuSei23}}]\label{thm:labrador-security}
    Let $\mathcal{C}$ be a challenge set as in \Cref{section:labrador-main} with parameters $\tau$ and $\theta$.
    Suppose that Module-SIS is hard for rank $\kappa_1 = \kappa_2$ and norm $2\beta'$, and for rank $\kappa$ and norm
    \begin{equation*}
        \max\Big\{8\theta(b+1)\beta', \; 2(b+1)\beta' + 4\theta\sqrt{128/30}\,\beta\Big\}.
    \end{equation*}
    Suppose further that $\beta \leq \sqrt{30/128}\cdot q/125$.
    Then the main protocol (\Cref{fig:labrador-main}) is a knowledge-sound proof-of-knowledge reduction for $\mathcal{R}_{\mathsf{dot}}$ with soundness error $2^{-125}$ and norm slack $\sqrt{128/30} \approx 2.07$, \textit{i.e.}, the extractor outputs a witness with $\sum_i \|\mathbf{s}_i\|_2^2 \leq \frac{128}{30}\beta^2$.
\end{theorem}

The theorem explains three concrete choices of LaBRADOR.
First, the condition $\beta \lesssim q/125$ comes from \Cref{lemma:labrador-jl}, and it allows a relatively small single-precision modulus $q \approx 2^{32}$ (in contrast, the version in \Cref{lemma:jl-projection} requires a gap between the norm and $q$ that grows with the dimension).
Second, the factor $\theta$ in the Module-SIS norm is the reason why the challenges are restricted to a small operator norm.
Third, the rank $\kappa$ of the inner commitments must be chosen large enough for the norm $\approx 8\theta(b+1)\beta'$, which grows with the decomposition base $b$; this is one of the trade-offs that determine the decomposition parameters.
\fi

\paragraph{Concrete parameters choices}
The implementation of LaBRADOR uses the ring $\mathcal{R}_q = \mathbb{Z}_q[T]/(T^{64}+1)$ (so $d = 64$) with a modulus $q \approx 2^{32}$ for which $T^{64}+1$ splits into two irreducible factors, the challenge set with $\tau = 71$ and $\theta = 15$ from \Cref{section:labrador-main}, and commitment ranks $\kappa$, $\kappa'$, $\kappa''$ chosen in each iteration so that Module-SIS offers $128$ bits of security.
The protocol is made non-interactive with the Fiat-Shamir transform, and all challenges (including the projection matrices) are expanded from short seeds.

\subsection{Reduction From R1CS to Dot Product Constraints}\label{section:labrador-r1cs}

It remains to connect LaBRADOR to the circuits of the previous chapters.
Recall the matrix form of R1CS: a witness $\mathbf{w}$ satisfies the system $(\mathbf{L}, \mathbf{R}, \mathbf{O})$ if
\begin{equation*}
    \mathbf{L}\mathbf{w} \odot \mathbf{R}\mathbf{w} = \mathbf{O}\mathbf{w},
\end{equation*}
where $\mathbf{L}, \mathbf{R}, \mathbf{O} \in \mathbb{Z}_N^{k \times \ell}$ describe $k$ constraints over $\ell$ variables.
LaBRADOR supports two moduli: $N = 2$ (\emph{binary} R1CS, convenient for hash functions and bit manipulations) and $N = 2^d + 1$ (convenient for integer arithmetic).
In both cases, the reduction is a short proof-of-knowledge reduction whose target relation is $\mathcal{R}_{\mathsf{dot}}$, and the main protocol then proves the target instance.

\subsubsection{Binary R1CS}

Let $\mathbf{w} \in \{0,1\}^{\ell}$ be a witness for a binary R1CS.
The prover computes the three vectors
\begin{equation*}
    \mathbf{a} := \mathbf{L}\mathbf{w} \bmod 2, \quad \mathbf{b} := \mathbf{R}\mathbf{w} \bmod 2, \quad \mathbf{c} := \mathbf{O}\mathbf{w} \bmod 2,
\end{equation*}
inside $\{0,1\}^k$, packs all bits of $\mathbf{a}$, $\mathbf{b}$, $\mathbf{c}$, and $\mathbf{w}$ into the coefficients of polynomials (padding with zeros so that $k$ and $\ell$ are multiples of $d$), and commits to them with an Ajtai commitment $\mathbf{t} = \mathbf{A}_{\mathsf{R1CS}}(\mathbf{a}\,\|\,\mathbf{b}\,\|\,\mathbf{c}\,\|\,\mathbf{w})$.
In what follows, we denote the packed vectors of polynomials by the same letters.
The prover then needs to show three things: all committed coefficients are bits, $\mathbf{a} \odot \mathbf{b} = \mathbf{c}$, and $\mathbf{a}$, $\mathbf{b}$, $\mathbf{c}$ are computed correctly from $\mathbf{w}$ modulo $2$.

\paragraph{Binarity}
Let $\mathbf{1}$ denote the vector of polynomials whose coefficients are all equal to $1$.
By \Cref{lemma:constant-coefficient},
\begin{equation*}
    \mathsf{ct}\big(\langle \mathbf{a}^{\ast}, \mathbf{a} - \mathbf{1} \rangle\big) = \sum_{i} a_i(a_i - 1) \bmod q,
\end{equation*}
where $a_i$ are the coefficients of $\mathbf{a}$.
Since $\mathbf{a}^{\ast}$ is not a dot product function of $\mathbf{a}$, the prover adds $\widetilde{\mathbf{a}} := \mathbf{a}^{\ast}$ to the witness of the target instance: the relation $\widetilde{\mathbf{a}} = \mathbf{a}^{\ast}$ is $\mathbb{Z}_q$-linear (conjugation only permutes the coefficients and flips some signs), so it becomes a set of constant-coefficient constraints in $\mathcal{F}_{\mathsf{ct}}$ (one per coefficient), and the constraint $\mathsf{ct}(\langle \widetilde{\mathbf{a}}, \mathbf{a} - \mathbf{1} \rangle) = 0$ belongs to the family $\mathcal{F}_{\mathsf{ct}}$.
This constraint holds modulo $q$, while we need it over the integers.
Here the norm bound of $\mathcal{R}_{\mathsf{dot}}$ comes into play: every term $a_i(a_i - 1)$ is a non-negative integer, and $\sum_i a_i(a_i - 1) \leq 2\|\mathbf{a}\|_2^2$.
Hence, if the target instance certifies that the squared norm of the witness (which contains both $\mathbf{a}$ and $\widetilde{\mathbf{a}}$, so it is at least $2\|\mathbf{a}\|_2^2$) is below $q$, then the sum cannot wrap around modulo $q$, it is zero over the integers, and all $a_i$ are bits.
The same argument applies to $\mathbf{b}$, $\mathbf{c}$, and $\mathbf{w}$.

\paragraph{Products}
For bits $a$, $b$, $c$, one has $ab = c$ if and only if $a + b - 2c \in \{0, 1\}$.
Hence, the prover shows that the coefficients of $\mathbf{a} + \mathbf{b} - 2\mathbf{c}$ are bits, using the same trick with the constraint $\mathsf{ct}(\langle \widetilde{\mathbf{a}} + \widetilde{\mathbf{b}} - 2\widetilde{\mathbf{c}}, \mathbf{a} + \mathbf{b} - 2\mathbf{c} - \mathbf{1} \rangle) = 0$.

\paragraph{Linear relations modulo $2$}
Finally, the prover must show that $\mathbf{a} = \mathbf{L}\mathbf{w} \bmod 2$, and similarly for $\mathbf{b}$ and $\mathbf{c}$.
These are $3k$ linear relations modulo $2$, while the main protocol proves relations modulo $q$.
The verifier first compresses them into $\lambda$ random combinations: it sends random $\boldsymbol{\alpha}_{i,L}, \boldsymbol{\alpha}_{i,R}, \boldsymbol{\alpha}_{i,O} \leftarrowS \{0,1\}^k$ for $i \in [\lambda]$, and both parties compute the binary vectors
\begin{equation*}
    \boldsymbol{\delta}_i := \boldsymbol{\alpha}_{i,L}^{\top}\mathbf{L} + \boldsymbol{\alpha}_{i,R}^{\top}\mathbf{R} + \boldsymbol{\alpha}_{i,O}^{\top}\mathbf{O} \bmod 2 \in \{0,1\}^{\ell}.
\end{equation*}
The prover sends the integers
\begin{equation*}
    e_i := \langle \boldsymbol{\alpha}_{i,L}, \mathbf{a} \rangle + \langle \boldsymbol{\alpha}_{i,R}, \mathbf{b} \rangle + \langle \boldsymbol{\alpha}_{i,O}, \mathbf{c} \rangle - \langle \boldsymbol{\delta}_i, \mathbf{w} \rangle, \quad i \in [\lambda],
\end{equation*}
where the dot products are computed over the integers on the bits of the vectors.
If all relations hold, then every $e_i$ is even, and the verifier checks this.
If some relation fails, then each $e_i$ is odd with probability $1/2$, so the soundness error is $2^{-\lambda}$.
The correctness of each $e_i$ is a $\mathbb{Z}_q$-linear relation on the coefficients of the witness, \textit{i.e.}, a constraint of the family $\mathcal{F}_{\mathsf{ct}}$; it holds over the integers as well, as long as $\ell + 3k < q$, since then $|e_i| < q$.

The complete reduction is specified in \Cref{fig:labrador-binary-r1cs}.
The target witness consists of the vectors $\mathbf{a}$, $\mathbf{b}$, $\mathbf{c}$, $\mathbf{w}$ and their conjugates, and its squared norm is at most $2(3k + \ell)$.
Hence, the condition $2(3k + \ell)\cdot\frac{128}{30} < q$, which holds for all practical circuit sizes with $q \approx 2^{32}$, leaves enough room for the norm slack of the main protocol.

\begin{algoproto}[H]
    \small
    \centering
    \begin{mdframed}[linecolor=black, linewidth=1pt,
                     skipabove=10pt, skipbelow=10pt,
                     innerleftmargin=5pt, innerrightmargin=5pt]
        \begin{minipage}[t]{0.5\linewidth}
            \colorbox{oc-blue-1}{$\mathcal{P}(\mathbbm{x} = (\mathbf{L}, \mathbf{R}, \mathbf{O}), \mathbbm{w} = \mathbf{w})$:}
            \begin{algoen}
                \item $\mathbf{a}, \mathbf{b}, \mathbf{c} \gets \mathbf{L}\mathbf{w}, \mathbf{R}\mathbf{w}, \mathbf{O}\mathbf{w} \bmod 2$;
                \item \keyword{send} $\mathbf{t} \gets \mathbf{A}_{\mathsf{R1CS}}(\mathbf{a}\,\|\,\mathbf{b}\,\|\,\mathbf{c}\,\|\,\mathbf{w})$;
                \item \keyword{receive} $(\boldsymbol{\alpha}_{i,L}, \boldsymbol{\alpha}_{i,R}, \boldsymbol{\alpha}_{i,O})_{i \in [\lambda]}$;
                \item compute $\boldsymbol{\delta}_i$ and $e_i$ for $i \in [\lambda]$;
                \item \keyword{send} $(e_i)_{i \in [\lambda]}$;
                \item \commentgray{target witness} \\ \keyword{return} $(\mathbf{a}, \mathbf{b}, \mathbf{c}, \mathbf{w}, \mathbf{a}^{\ast}, \mathbf{b}^{\ast}, \mathbf{c}^{\ast}, \mathbf{w}^{\ast})$.
            \end{algoen}
        \end{minipage}
        \hfill
        \begin{minipage}[t]{0.46\linewidth}
            \colorbox{oc-blue-1}{$\mathcal{V}(\mathbbm{x} = (\mathbf{L}, \mathbf{R}, \mathbf{O}))$:}
            \begin{algoen}
                \item \keyword{receive} $\mathbf{t}$;
                \item \keyword{send} $\boldsymbol{\alpha}_{i,L}, \boldsymbol{\alpha}_{i,R}, \boldsymbol{\alpha}_{i,O} \leftarrowS \{0,1\}^k$ for $i \in [\lambda]$;
                \item \keyword{receive} $(e_i)_{i \in [\lambda]}$;
                \item \keyword{assert} $e_i$ is even for all $i \in [\lambda]$;
                \item \commentgray{target statement} \\ \keyword{return} $(\mathcal{F}, \mathcal{F}_{\mathsf{ct}}, \sqrt{2(3k+\ell)})$.
            \end{algoen}
        \end{minipage}
    \end{mdframed}
    \caption{Proof-of-knowledge reduction from binary R1CS to $\mathcal{R}_{\mathsf{dot}}$~\cite[Fig.~4]{C:BeuSei23}. The family $\mathcal{F}$ contains the opening $\mathbf{t} = \mathbf{A}_{\mathsf{R1CS}}(\mathbf{a}\,\|\,\mathbf{b}\,\|\,\mathbf{c}\,\|\,\mathbf{w})$, while the family $\mathcal{F}_{\mathsf{ct}}$ contains the constant-coefficient constraints for the conjugation relations such as $\widetilde{\mathbf{a}} = \mathbf{a}^{\ast}$, for the binarity of $\mathbf{a}$, $\mathbf{b}$, $\mathbf{c}$, $\mathbf{w}$, and $\mathbf{a} + \mathbf{b} - 2\mathbf{c}$, and for the correctness of the $e_i$.}
    \label{fig:labrador-binary-r1cs}
\end{algoproto}

\iffalse

\subsubsection{R1CS Modulo $2^d + 1$}

Binary circuits are inconvenient for arithmetic. For example, a single multiplication of two $64$-bit integers requires thousands of binary constraints.
For such statements, LaBRADOR uses R1CS modulo $2^d + 1$, and the reason for this particular modulus is the following ring homomorphism.
Since $2^d \equiv -1 \pmod{2^d + 1}$, the number $2$ is a root of $T^d + 1$ modulo $2^d + 1$, so evaluation at $T = 2$ induces the ring homomorphism
\begin{equation*}
    \mathsf{ev}: \mathcal{R} \to \mathbb{Z}_{2^d+1}, \quad \sum_{i < d} f_iT^i \mapsto \sum_{i < d} f_i2^i \bmod (2^d + 1).
\end{equation*}

\begin{remark}
    This is the same observation as in Exercise~11 of \Cref{section:lattices-exercises}.    
\end{remark}

Now we observe the following: every $x \in \mathbb{Z}_{2^d+1}$ has a short preimage under $\mathsf{ev}$. 
Indeed, writing $x$ in the \emph{non-adjacent form}\footnote{\emph{Non-adjacent decomposition} is a decomposition in base $2$ with digits in $\{-1, 0, 1\}$ where no two adjacent digits are non-zero.} gives a polynomial $\mathsf{Enc}(x) \in \mathcal{R}$ with coefficients in $\{-1, 0, 1\}$ and $\|\mathsf{Enc}(x)\|_2 \leq \sqrt{d/2}$.
Products of encodings are encodings of products: $\mathsf{ev}(\mathsf{Enc}(x)\cdot\mathsf{Enc}(y)) = xy$, as long as the product computed in $\mathcal{R}_q$ does not wrap around modulo $q$.

\begin{exercise}\label{exercise:labrador-ev}
    Let $d = 4$, so that $2^d + 1 = 17$.
    Find the non-adjacent forms $\mathsf{Enc}(3)$ and $\mathsf{Enc}(5)$, compute their product in $\mathbb{Z}[T]/(T^4 + 1)$, and check that $\mathsf{ev}$ maps it to $3 \cdot 5 \bmod 17$.

    \commentgray{Answer: $\mathsf{Enc}(3) = T^2 - 1$, $\mathsf{Enc}(5) = T^2 + 1$, and their product is $T^4 - 1 = -2$, which evaluates to $-2 \equiv 15 = 3 \cdot 5 \pmod{17}$.}
\end{exercise}

With this encoding, the prover commits to the encodings of $\mathbf{L}\mathbf{w}$, $\mathbf{R}\mathbf{w}$, $\mathbf{O}\mathbf{w}$, and $\mathbf{w}$, and proves that they are well-formed and that $\mathbf{L}\mathbf{w} \odot \mathbf{R}\mathbf{w} = \mathbf{O}\mathbf{w}$.
As in the binary case, the linear relations are compressed with random combinations, and the $k$ products are compressed into $\lambda$ random inner products $\langle \boldsymbol{\phi}_i, \mathbf{L}\mathbf{w} \odot \mathbf{R}\mathbf{w} - \mathbf{O}\mathbf{w} \rangle = 0$ for $\boldsymbol{\phi}_i \leftarrowS \mathbb{Z}_{2^d+1}^k$.
Each resulting relation is a dot product function of the encodings, and the prover sends its value $e_i \in \mathcal{R}_q$.
The verifier lifts $e_i$ to a polynomial with integer coefficients and checks that $\mathsf{ev}(e_i) = 0$, \textit{i.e.}, that $e_i$ is divisible by $T - 2$; the lift is exact because the parameters are chosen so that no wraparound modulo $q$ occurs.
Since $2^d + 1$ is not prime, the soundness error of a random combination is $1/p$, where $p$ is the smallest prime factor of $2^d + 1$.
For $d = 64$, one has $2^{64} + 1 = 274177 \cdot 67280421310721$, so $p \approx 2^{18}$ and $\lambda = \lceil 128/18 \rceil = 8$ combinations suffice.

\paragraph{Mixing both types of constraints}
Many statements need both kinds of constraints.
For example, proving knowledge of a Dilithium signature requires a lot of arithmetic modulo $8380417$, which is expensive as a binary circuit, but also norm checks and hashing, which are expensive as an arithmetic circuit.
Since both reductions produce instances of $\mathcal{R}_{\mathsf{dot}}$, the prover can commit to a binary witness that simultaneously encodes the witness of the arithmetic R1CS, run both reductions in parallel, and prove the union of the two target instances with a single execution of the recursive protocol.
\fi

\end{document}